% Figure: recursive (online) estimation. The running estimate of a
% constant converges toward the truth as measurements arrive, and the
% uncertainty band shrinks as 1/sqrt(n).
\begin{figure}[htbp]
\centering
\begin{tikzpicture}[font=\small]
\draw[->, thick, draw=engblack!70] (0,0) -- (9.2,0)
  node[right, font=\scriptsize] {measurement number $n$};
\draw[->, thick, draw=engblack!70] (0,0) -- (0,4.6)
  node[above, font=\scriptsize] {estimate};
% truth line
\draw[thick, draw=enggray, dashed] (0,2.6) -- (8.8,2.6);
\node[font=\scriptsize, color=enggray, anchor=west] at (8.0,2.8) {truth};
% uncertainty band: shrinks like 1/sqrt(n) around the running mean.
% Build upper and lower envelopes as piecewise lines.
% running mean wanders then settles at 2.6; band half-width ~ 1.6/sqrt(n)
\foreach \n/\m in {1/3.7, 2/3.2, 3/2.2, 4/2.9, 5/2.45, 6/2.75, 7/2.5, 8/2.62, 9/2.58, 10/2.6} {
  \pgfmathsetmacro\xx{0.85*\n}
  \pgfmathsetmacro\hw{1.7/sqrt(\n)}
}
% band as a filled region (approximate): sample the envelope
\filldraw[engblue!12, draw=none]
  (0.85,{3.55+1.7}) --
  (1.70,{3.10+1.7/sqrt(2)}) --
  (2.55,{2.70+1.7/sqrt(3)}) --
  (3.40,{2.75+1.7/sqrt(4)}) --
  (4.25,{2.60+1.7/sqrt(5)}) --
  (5.10,{2.68+1.7/sqrt(6)}) --
  (5.95,{2.58+1.7/sqrt(7)}) --
  (6.80,{2.62+1.7/sqrt(8)}) --
  (7.65,{2.60+1.7/sqrt(9)}) --
  (8.50,{2.60+1.7/sqrt(10)}) --
  (8.50,{2.60-1.7/sqrt(10)}) --
  (7.65,{2.60-1.7/sqrt(9)}) --
  (6.80,{2.62-1.7/sqrt(8)}) --
  (5.95,{2.58-1.7/sqrt(7)}) --
  (5.10,{2.68-1.7/sqrt(6)}) --
  (4.25,{2.60-1.7/sqrt(5)}) --
  (3.40,{2.75-1.7/sqrt(4)}) --
  (2.55,{2.70-1.7/sqrt(3)}) --
  (1.70,{3.10-1.7/sqrt(2)}) --
  (0.85,{3.55-1.7}) -- cycle;
% running estimate line
\draw[very thick, draw=engblue!85]
  (0.85,3.55) -- (1.70,3.10) -- (2.55,2.70) -- (3.40,2.75) --
  (4.25,2.60) -- (5.10,2.68) -- (5.95,2.58) -- (6.80,2.62) --
  (7.65,2.60) -- (8.50,2.60);
% individual noisy measurements
\foreach \xx/\yy in {0.85/4.3, 1.70/2.1, 2.55/1.6, 3.40/3.4, 4.25/2.0, 5.10/3.1, 5.95/2.1, 6.80/2.9, 7.65/2.4, 8.50/2.7} {
  \filldraw[engred!70] (\xx,\yy) circle (0.05);
}
\node[font=\scriptsize, color=engred!70, anchor=west] at (1.9,1.4)
  {individual measurements};
\node[font=\scriptsize, color=engblue!85, anchor=south west] at (3.6,3.0)
  {running estimate $\pm$ band};
\end{tikzpicture}
\caption[Recursive estimation of a constant]{A running estimate of a
constant quantity updated as each measurement arrives. The individual
measurements scatter widely; the recursive estimate, which blends the
prior estimate with each new datum in proportion to their relative
precisions, converges toward the truth. The shaded uncertainty band
narrows as $1/\sqrt{n}$, the same rate the batch mean achieves, but the
recursive form produces the current best estimate and its interval at
every step without re-processing the whole history. This is the update
rule at the heart of online identification and of the Kalman filter the
later volumes develop; the modelling content is that an estimate carries
its uncertainty forward and shrinks it by exactly the information each new
measurement adds.}
\label{fig:vol02ch18:recursive-estimation}
\end{figure}
